\documentclass[10pt,a4paper]{amsart}
\usepackage[margin=19mm]{geometry}
\usepackage{amsmath,amssymb,mathtools}
\usepackage[scr=rsfs,cal=euler]{mathalfa}
\usepackage{xfrac}
\usepackage[unicode,hidelinks,bookmarks=false]{hyperref}
\newcommand{\ie}{i.e.\ }
\newcommand{\cf}{cf.\ }
\newcommand{\resp}{respectively\ }
\newcommand{\Subsection}{\subsection}
\providecommand{\nobreakdash}{\nobreak\hbox{-}}
\setlength{\emergencystretch}{2em}
\multlinegap0pt
\usepackage{amssymb}
\usepackage[shortlabels]{enumitem}
\usepackage{tikz}
\newtheorem{theorem}{Theorem}
\newtheorem{corollary}{Corollary}
\usepackage{cleveref}
\crefname{theorem}{Theorem}{Theorems}
\crefname{corollary}{Corollary}{Corollarys}
\DeclareMathOperator{\Gr}{Gr}
\title{Symbolic Powers and Asymptotic Invariants of GL-Invariant Ideals}
\author{Sankhaneel Bisui, Alexandra Seceleanu}
\date{Source: arXiv:2609.03131v1 (CC BY 4.0)}
\begin{document}
\maketitle

\noindent\emph{Source and licence.} Symbolic Powers and Asymptotic Invariants of GL-Invariant Ideals. Version arXiv:2609.03131v1 (CC BY 4.0), distributed by arXiv under the Creative Commons Attribution 4.0 International licence, http://creativecommons.org/licenses/by/4.0/, as recorded in the arXiv licence metadata for this exact version and shown on the abstract page https://arxiv.org/abs/2609.03131v1. Full licence notice: \texttt{licenses/arxiv-2609.03131-CC-BY-4.0.txt}.\par\smallskip

\noindent\textit{Editorial scope.} Counted unit: the article's corollary corrected (source0.tex lines 1845--1854). The article's complete original proof of it is delivered with it (source0.tex lines 1856--1865, one \texttt{proof} environment of 10 lines). No mathematics is written, repaired or re-derived: the statement and the proof are the article's own lines, in the article's own notation, and no retained line is substituted or re-flowed.  Retained uncounted as the source-local context the proof needs: lines 1723--1729.  Nothing was omitted from any retained span, so the package carries no omission record.  No retained line contains a citation, so the wrapper carries no bibliography and no bibliography entry.  No retained line contains a figure environment, an image-inclusion command or drawing code, so no image is delivered and the package declares no asset dependency.  Editorial formatting only, in addition: an AMS \texttt{amsart} preamble that restates the article's own declarations and macros used by the retained lines, namely the environments \texttt{theorem}, \texttt{corollary} on separate counters, together with the standard packages those lines need.  The retained environments print in this document as Theorem 1, Corollary 1 in order of appearance, whereas the article prints the same items with the numbering of its own full document.  The complete native source of the article as distributed in the arXiv e-print for this exact version is bundled unchanged as \texttt{sources/209267/source0.tex}.
\begin{theorem}\label{thm: rectangle corrected}
Let $\sigma=(\sigma_1,\ldots,\sigma_s)\in H_r$, let
$\ell=|\sigma|$, and set $\rho=(r^\ell)$.  Then
\[
c_{\underbrace{\scriptstyle\sigma^\vee,\ldots,\sigma^\vee}_{r\text{ copies}}}^{\rho^\vee}>0.
\]
\end{theorem}

\begin{corollary}\label{cor: Schubert corrected}
Let $k,r$ be positive integers, and let $\mathcal S_\sigma$ be the Schubert
class indexed by a partition $\sigma$ satisfying $|\sigma|\leq k$ and
$\sigma_1\leq r$ in $H^*(\Gr(k,k+r),\mathbb Z)$.  Then
$\mathcal S_\sigma^r\neq0$.  If $|\sigma|=k$, then
\[
\mathcal S_\sigma^r=c[\mathrm{pt}]
\qquad\text{for some integer }c>0.
\]
\end{corollary}

\begin{proof}
Set $\ell=|\sigma|$.  The rectangle $(r^\ell)$ fits inside the $k\times r$
rectangle because $\ell\leq k$.  By \Cref{thm: rectangle corrected}, its
coefficient in $\mathcal S_\sigma^r$ is positive, proving nonvanishing.  If
$\ell=k$, this rectangle is the full $k\times r$ rectangle and represents the
point class.  Since the product already has top degree, the displayed formula
follows.  The last assertion records the distinction between a nonzero
multiple of the generator of the rank-one group
$H^{2kr}(\Gr(k,k+r),\mathbb Z)$ and an integral generator.
\end{proof}

\end{document}
